\section{Onde aprofundar}

\subsection*{Papers}
\small
\begin{longtable}{@{}p{5.2cm} p{7cm}@{}}
\toprule
Fonte & Para quê \\
\midrule \endhead
\textcite{paper:statistically-rigorous-java-perf} & Leitura obrigatória
  antes de comparar dois números de JVM \\
\textcite{paper:vm-warmup-hot-cold} & Por que \enquote{basta aquecer} nem
  sempre é verdade; versão aberta no arXiv \\
\textcite{paper:dacapo} e \textcite{paper:dacapo-chopin} & Como uma suíte
  de aplicações reais é montada e analisada \\
\textcite{paper:renaissance} & Concorrência e programação funcional como
  carga de benchmark; PDF aberto no site da suíte \\
\bottomrule
\end{longtable}
\normalsize

\subsection*{Livros}
\citetitle{livro:oaks2020} e \citetitle{livro:evans2018} tratam de
medição de desempenho na JVM do ponto de vista de quem opera serviços.

\subsection*{Ferramentas}
\citetitle{doc:jmh} para microbenchmarks; \citetitle{doc:renaissance} e
\citetitle{doc:dacapo} para cargas de aplicação; o guia
\citetitle{doc:quarkus-performance-measure} e o repositório
\citetitle{doc:quarkus-perf-comparison} para medir um serviço inteiro.
